% Ce chapitre repose sur :
%   DOC: (aucun)
%   OBS: REL-RDV.C02, REL-RDV.R08
%   HYP: (aucun)
%   CHI: agregats, colonnes-intermediate-typees, dim-patient-controles-versions, dim-patient-fiches-plusieurs-versions, dim-patient-lignes, dimensions
%   CHI: dimensions-sans-libelle, faits, identifiants-distincts, instantane-disparition-vues, modeles-intermediate, modeles-marts
%   CHI: modeles-total, partitions-par-table, plancher-rejets, quarantaine-lignes-total, quarantaine-patients-lignes, quarantaine-rendez-vous-lignes
%   CHI: seuil-quarantaine, source-colonnes-texte, source-lignes-total, source-tables, taches-graphe
%   SER: (aucune)
%   SECTIONS: 6
\chapter{Architecture de la chaîne de données}\label{chap:architecture}

La chaîne transforme les fichiers extraits en tables qu'un tableau de bord peut lire. La question
traitée est celle de son découpage : quelles couches elle traverse, pourquoi chacune existe, et
comment elles s'enchaînent chaque jour.

\section{Vue d'ensemble}

La chaîne va d'un dépôt de fichiers à un tableau de bord en cinq étapes (voir
figure~\ref{fig:architecture}).

\begin{figure}[htbp]
\centering
\begin{tabular}{@{}c@{}}
\fbox{\begin{minipage}{0.90\linewidth}\centering\small
\fbox{\begin{minipage}{0.94\linewidth}\centering \textbf{Zone d'atterrissage} \\
\footnotesize fichiers partitionnés par table et par date d'extraction\end{minipage}}\\[2pt]
$\downarrow$ \footnotesize chargement : contrôles d'entrée, quarantaine, transaction par partition\\[2pt]
\fbox{\begin{minipage}{0.44\linewidth}\centering \texttt{source} \\
\footnotesize copie fidèle, tout en texte\end{minipage}}%
\hfill
\fbox{\begin{minipage}{0.44\linewidth}\centering \texttt{quarantaine} \\
\footnotesize lignes rejetées, avec leur motif\end{minipage}}\\[2pt]
$\downarrow$ \footnotesize transformation : typage, normalisation\\[2pt]
\fbox{\begin{minipage}{0.94\linewidth}\centering \texttt{intermediate} \\
\footnotesize une vue par table source, colonnes typées\end{minipage}}\\[2pt]
$\downarrow$ \footnotesize modélisation dimensionnelle\\[2pt]
\fbox{\begin{minipage}{0.94\linewidth}\centering \texttt{marts} \\
\footnotesize dimensions, faits, agrégats\end{minipage}}\\[2pt]
$\downarrow$ \footnotesize rapprochement d'identités \quad$\rightarrow$\quad \texttt{linkage}\\[2pt]
$\downarrow$ \footnotesize échange de noms en une transaction\\[2pt]
\fbox{\begin{minipage}{0.94\linewidth}\centering \texttt{instantane} \\
\footnotesize tables figées que le tableau de bord lit\end{minipage}}
\end{minipage}}
\end{tabular}
\caption{Architecture de la chaîne de données}\label{fig:architecture}
\end{figure}

Aucune couche ne modifie la précédente. La couche source copie les fichiers avec leurs défauts, la
couche intermédiaire type et normalise sans juger, et la couche analytique modélise. Corriger à
l'entrée effacerait ce qu'il faut mesurer, et corriger en sortie empêcherait de rejouer la chaîne.

Les outils de chaque couche sont présentés en section~\ref{sec:complements-pile}, et le
chapitre~\ref{chap:outils} les nomme. La chaîne compte \chiffre{modeles-total} modèles de transformation, dont
\chiffre{modeles-intermediate} intermédiaires et \chiffre{modeles-marts} analytiques, et
\chiffre{taches-graphe} tâches d'orchestration. Le tableau de bord ne lit pas la couche analytique
mais un instantané, copie figée de l'état construit, décrit en section~\ref{sec:orchestration}.

\section{Zone d'atterrissage et couche source}

La zone d'atterrissage range les fichiers par table et par date d'extraction, et la couche source
les copie sans contrainte. Chaque table compte \chiffre{partitions-par-table} partitions sur la
période, une par jour. Ce découpage est une construction de la chaîne, absente du système observé.
Il rend le chargement rejouable jour par jour et donne à chaque ligne une date d'extraction, qui
devient une colonne.

Chaque partition est écrite en deux formats : un fichier de valeurs séparées, que le chargeur lit,
et un fichier colonnaire livré pour l'échange. Un contrôle vérifie que les deux portent le même
contenu, colonne par colonne. Chaque ligne doit aussi déclarer la date d'extraction du
répertoire qui la contient, sans quoi elle part en quarantaine.

Les \chiffre{source-tables} tables de la couche source comptent \chiffre{source-colonnes-texte}
colonnes, toutes de type texte, sans clé ni contrainte. Une contrainte ferait échouer le chargement
sur une ligne défectueuse, alors que ces lignes sont la matière du travail de qualité. Le typage a
lieu à la couche suivante, où il devient une transformation observable. Sur la période, la couche
source porte \chiffre{source-lignes-total} lignes.

\section{Le choix du serveur de base de données}\label{sec:choix-serveur}

Le serveur relationnel a été retenu pour trois propriétés que la chaîne utilise vraiment. Plusieurs
services interrogent la même base par le réseau, avec des identifiants distincts. Le chargement
écrit pendant que le tableau de bord lit, ce qui demande des transactions concurrentes. Le catalogue
de la base sert enfin à compter les colonnes que le registre déclare. Un moteur embarqué aurait
produit les mêmes tables, et plus vite sur ce volume : le choix ne tient donc pas à la performance.

Le serveur tourne dans un conteneur, si bien que la base est la même sur le poste de développement
et dans l'intégration continue. Seul le rapprochement d'identités s'exécute sur un moteur en
mémoire. Il écrit de nombreuses tables intermédiaires, et les tenir hors du serveur évite d'y
laisser des résidus.

\section{Couche intermédiaire et quarantaine}\label{sec:quarantaine}

La couche intermédiaire compte \chiffre{modeles-intermediate} modèles, un par table source : elle
type, normalise et ne rejette rien. Les formats de date ont été mesurés avant d'être écrits. Au
poste, les dates sont au format mois-jour-année, et les horodatages y ajoutent l'indicateur du
matin ou de l'après-midi. Ce format a été relevé sur la date de création d'un
rendez-vous\releve{REL-RDV.C02} comme sur sa date de tenue\releve{REL-RDV.R08}.

La conversion est écrite une fois, dans une macro, et appliquée partout. Sans elle, le jour et le
mois s'échangeraient sans bruit onze fois sur douze. Les mesures faites sur la base avant d'écrire
les conversions ont aussi montré trois faits : la valeur manquante est toujours la chaîne vide et
jamais l'absence de valeur, les dates et les horodatages se convertissent tous sur les
\chiffre{colonnes-intermediate-typees} colonnes typées, et le domaine booléen compte exactement deux
valeurs et la chaîne vide.

Une ligne refusée au chargement est mise en quarantaine : elle est écrite dans un schéma parallèle,
avec le motif de son rejet et sa partition, et jamais supprimée. Sur la période,
\chiffre{quarantaine-lignes-total} lignes y sont, dont \chiffre{quarantaine-patients-lignes} fiches
patient et \chiffre{quarantaine-rendez-vous-lignes} rendez-vous, pour
\chiffre{source-lignes-total} lignes chargées. Ce sont les défauts injectés que les contrôles
d'entrée ont attrapés : dates hors domaine, âges incohérents avec la date de naissance, valeurs hors
nomenclature. Les autres tables n'en portent aucune, parce que leurs défauts ne se voient pas sur
une ligne seule, et garder ces lignes permet de mesurer la qualité d'entrée.

Un seuil bloque le chargement d'un fichier trop défectueux. Au-delà de
\chiffre{seuil-quarantaine} de lignes rejetées, le fichier est bloqué plutôt que chargé en partie,
car un tel taux signale plutôt un défaut d'extraction qu'une suite de défauts de saisie. Un plancher
de \chiffre{plancher-rejets} rejets évite qu'un seul rejet, dans une petite partition, dépasse le
seuil à lui seul, et il faut les deux conditions pour bloquer. Un second contrôle compare au même
seuil le cumul de la journée, toutes tables confondues, pour qu'une mauvaise journée ne passe pas
fichier par fichier.

\section{Modélisation dimensionnelle}

La couche analytique suit un schéma en étoile : chaque table de faits porte un événement, et les
dimensions qui l'entourent le décrivent (voir figure~\ref{fig:schema-etoile}). Elle compte
\chiffre{modeles-marts} modèles, qui produisent \chiffre{dimensions} dimensions,
\chiffre{faits} faits et \chiffre{agregats} agrégats. Les deux décomptes diffèrent parce qu'un des
agrégats est une vue posée hors de l'outil de transformation (section~\ref{sec:complements-analytique}).

Sur les \chiffre{dimensions} dimensions, \chiffre{dimensions-sans-libelle} ne portent que leur code, sans libellé. Aucune source ne nomme ces
libellés, et en écrire un plausible afficherait une hypothèse sous l'apparence d'une donnée
relevée. Le calendrier est construit, et la dimension patient garde l'historique des fiches.

Les agrégats portent des décomptes déjà calculés pour l'affichage, comme l'activité par jour ou
l'absentéisme par activité. Ils existent pour que le tableau de bord rende une page en un temps
borné. Le registre des indicateurs dit, pour chaque valeur affichée, de quelles tables elle est
recalculée. Les créances et les relances n'ont pas de table de faits, et leurs agrégats partent de
la couche intermédiaire.

La dimension patient est la seule à garder une ligne par version de fiche. Un passage de janvier
reste ainsi rattaché à l'adresse connue en janvier, même si elle a été corrigée en mai. Elle porte
\chiffre{dim-patient-lignes} lignes pour \chiffre{identifiants-distincts} identifiants, dont
\chiffre{dim-patient-fiches-plusieurs-versions} ont plus d'une version. La clé des versions et les \chiffre{dim-patient-controles-versions} contrôles qui la vérifient sont décrits en section~\ref{sec:complements-analytique}.

\section{Orchestration, idempotence et rattrapage}\label{sec:orchestration}

Un orchestrateur enchaîne chaque jour les \chiffre{taches-graphe} tâches de la chaîne, dans l'ordre
que montre la figure~\ref{fig:sequence-quotidienne}.

\begin{figure}[htbp]
\centering
\begin{tikzpicture}[x=1mm, y=1mm]

% --- les six participants, et leurs lignes de vie ------------------------------
\foreach \i/\x/\nom in {1/0/{La zone\\d'atterrissage}, 2/23/{Le chargeur}, 3/46/{La base},
                        4/69/{L'outil de\\transformation}, 5/92/{Le rapprochement},
                        6/115/{L'instantané}}{
  \node[partUML] (p\i) at (\x,0) {\nom};
  \draw[vieUML] (p\i.south) -- ++(0,-126);
}

\def\msg#1#2#3#4{\draw[msgUML] (#1.south|-0,#3) -- (#2.south|-0,#3) node[etiqUML,midway] {#4};}

\msg{p1}{p2}{-8}{1. vérifier la disponibilité de l'extraction}
\msg{p2}{p3}{-19}{2. charger la journée}

% Le contrôle de qualité s'exécute sur la base elle-même : message réflexif.
\draw[msgUML] (p3.south|-0,-28) -- ++(7,0) -- ++(0,-5) -- ++(-7,0);
\node[etiqUML, anchor=west] at (54,-30) {3. contrôle de qualité — taux de rejet cumulé};

\msg{p3}{p4}{-42}{4. couche intermédiaire}
\msg{p4}{p3}{-50}{5. \textbf{les dimensions}}
\msg{p3}{p4}{-58}{6. \textbf{puis les faits}}
\msg{p4}{p3}{-66}{7. tests de la couche}
\msg{p3}{p5}{-76}{8. prédiction des paires}
\msg{p5}{p3}{-84}{9. regroupement et évaluation}
\msg{p3}{p4}{-94}{10. agrégats}
\msg{p3}{p6}{-104}{11. rafraîchir l'instantané}

% Le rafraîchissement échange les noms des tables : réflexif, et en une seule transaction.
\draw[msgUML] (p6.south|-0,-110) -- ++(7,0) -- ++(0,-5) -- ++(-7,0);
\node[etiqUML, anchor=east] at (113,-112) {\textbf{échange de tous les noms en une seule transaction}};

% Le douzième message sort du diagramme : il n'a pas de destinataire parmi les participants.
\draw[msgUML] (p3.south|-0,-122) -- ++(-40,0) node[etiqUML,midway] {12. exporter — fichiers de livraison} -- ++(-12,0);
\fill[accentgarde] (p3.south|-0,-122) ++(-52,0) circle (0.6mm);
\end{tikzpicture}
\caption{Le traitement quotidien entre les participants de la chaîne}\label{fig:sequence-quotidienne}
\end{figure}

Les dimensions se construisent avant les faits, parce qu'un fait se rattache à une version de
dimension. La construction des modèles tourne sur un seul fil : en parallèle, un modèle lirait un
objet en cours de remplacement. Les identifiants des tâches dans le code sont listés en
section~\ref{sec:complements-analytique}.

Le chargement est idempotent : le refaire donne le même état que le faire une fois. Le fichier du
12~mars est par exemple chargé, puis trouvé incomplet, corrigé et rechargé. Le chargement supprime
d'abord, dans la même transaction, tout ce qui porte la date du 12~mars dans la couche source et en
quarantaine, puis il insère. La journée est remplacée, jamais dédoublée.

Le rattrapage en découle : une journée manquante se charge sans toucher aux autres. Si les journées
du 12~mars, du 13~mars et du 14~mars manquent, les charger dans n'importe quel ordre donne le même
état final. Deux contrôles le vérifient, l'un en chargeant deux fois la même partition, l'autre en
chargeant les mêmes partitions dans deux ordres différents.

La production et l'application du schéma restent hors du graphe. Leurs fichiers de définition
suppriment les tables avant de les recréer, et les placer dans un graphe quotidien viderait la base
chaque nuit. Elles s'exécutent à la main, quand le modèle change.

Le graphe finit par deux tâches de sortie. La première rafraîchit l'instantané que lit le tableau
de bord, fait de tables et non de vues. Une reconstruction de la couche analytique fait en effet
disparaître ses vues pendant environ \chiffre{instantane-disparition-vues} seconde, et un lecteur y rencontre alors une
erreur d'objet inexistant. Le rafraîchissement bâtit donc des tables neuves sous des noms
provisoires, puis échange tous les noms en une seule transaction.

Le rafraîchissement copie chaque vue analytique et chaque table de rapprochement, plus deux tables
de service. Un lecteur ne devrait jamais voir un instantané à moitié neuf, mais aucun contrôle du
dépôt ne le vérifie. La seconde tâche exporte un fichier par table, régénéré à chaque exécution et
non suivi par le gestionnaire de versions.

\section*{Synthèse}

La chaîne copie, type, modélise puis fige les données, et aucune étape ne corrige la précédente.
Les lignes refusées sont gardées en quarantaine avec leur motif. La couche analytique suit un schéma
en étoile, et seule la dimension patient garde l'historique des fiches. Le chargement se rejoue jour
par jour sans rien dédoubler, et le tableau de bord lit un instantané échangé en une transaction.
Le chapitre~\ref{chap:qualite} mesure la qualité des données que cette chaîne produit.
